% 天体物理学（综述）
% license CCBYSA3
% type Wiki

本文根据 CC-BY-SA 协议转载翻译自维基百科\href{https://en.wikipedia.org/wiki/Astrophysics?utm_source=chatgpt.com}{相关文章}。


天体物理学是一门利用物理学与化学的方法和原理来研究天体及其现象的科学。\(^\text{[2]}\)作为该学科的奠基人之一，詹姆斯·基勒（James Keeler）曾指出，天体物理学“旨在探明天体的本质，而非它们在空间中的位置或运动——研究它们是什么，而非它们在哪里”，\(^\text{[3]}\)这一区别正体现于天体力学的研究范畴中。

天体物理学的研究对象包括太阳（太阳物理学）、其他恒星、星系、系外行星、星际介质以及宇宙微波背景辐射。\(^\text{[4][5]}\)科研人员在整个电磁波谱范围内观测这些天体的辐射，分析的物理特性包括光度、密度、温度及化学成分等。由于天体物理学的研究范围极其广泛，天体物理学家在研究中会综合运用多个物理学分支的概念与方法，包括经典力学、电磁学、统计力学、热力学、量子力学、相对论、核物理与粒子物理，以及原子与分子物理等。

在实践中，现代天文学研究通常涉及大量理论与观测物理学工作。天体物理学的主要研究领域包括：暗物质、暗能量、黑洞及其他天体的性质，以及宇宙的起源与最终命运。\(^\text{[4]}\)
理论天体物理学还研究多个主题，例如：太阳系的形成与演化、恒星动力学与演化、星系的形成与演化、磁流体力学、宇宙大尺度物质结构的形成、宇宙射线的起源、广义相对论与狭义相对论、量子宇宙学与物理宇宙学（即对宇宙最大尺度结构的物理研究），其中也包括弦论宇宙学与天体粒子物理等内容。
